\section{Future Work}
\label{sec:future_work}

While AgentSQL proves effective, the evolution of Text-to-SQL parsing demands deeper resilience and scalable strategies. Future iterations of this work will explore several key directions. First, to improve resilience to benchmark anomalies, since recent analyses reveal that prevalent benchmarks like BIRD contain significant annotation errors \cite{floratou2024broken}, future frameworks must integrate confidence-scoring or probabilistic validation to identify ambiguous or incorrect ground-truth schemas dynamically, moving beyond brittle deterministic evaluation. Second, regarding execution-feedback enhancement, while our Reflector--Critic loop currently isolates logical and syntactical flaws, systems like ReViSQL \cite{zhu2026revisqlachievinghumanleveltexttosql} and AGENTIQL \cite{heidari_agentiql_2025} have demonstrated that iterative execution feedback---directly observing database runtime states---is crucial for resolving deep semantic discrepancies; hence, we plan to integrate execution-driven empirical feedback loops to refine the Critic's correction phase. Third, to enable orchestrated test-time scaling to push performance toward human-expert levels, we aim to incorporate strategies such as diverse synthesis, iterative refinement, and tournament selection, drawing inspiration from Agentar-Scale-SQL \cite{wang_agentar-scale-sql_2025}, which can dramatically improve candidate SQL recall on extremely challenging queries. Fourth, we intend to work on MAGIC checklist expansion by extending the correction guidelines with patterns for window functions, such as \texttt{ROW\_NUMBER} or \texttt{LEAD/LAG}, and Common Table Expressions to address the specific failure modes identified in our experiments. Finally, we plan to implement streaming index updates by adapting the offline FAISS indexing mechanism for incremental schema updates, enabling real-time deployment on databases with rapidly evolving schemas.
